\chapter*{不定积分每日一题}

\refer{《积分的方法与技巧》金玉明\ 顾新身\ 毛瑞庭\ /\ 编著}

\exercise{求积分$\int{\tan^n{x}\d x}$}%在花括号内输入题目

\teacher{张晓鹏}%在花括号内输入你的名字

\postdate{2022年12月31日}%在花括号内输入你想要发布那一天的日期

\bigresult{
    \formula{\int{\tan^{2m+1}{x}\d x}=\sum_{r=0}^{m-1}\qty(-1)^{r+2}\frac{\tan^{2m-2r}{x}}{2m-2r}+\qty(-1)^m\ln{\abs{\sec{x}}}+C}
    \formula{\int{\tan^{2m}{x}\d x}=\sum_{r=0}^{m-1}\qty(-1)^r\frac{\tan^{2m-2r-1}{x}}{2m-2r-1}+\qty(-1)^mx+C}
}%在花括号内输入本题结果

\analysis{
    前面我们探索了$\sec{x}$和$\csc{x}$，今天我们来探索$\tan{x}$的高次积分，求解方法主要借助的是一个积分递推公式，下面直接看吧.
    \formula{
        \begin{aligned}
            \int{\tan^{x}\d x}&=\int{\tan^{n-2}{x}\qty(\sec^2{x}-1)\d x}\\
            &=\int{\tan^{n-2}{x}\sec^2{x}\d x}-\int{\tan^{n-2}{x}\d x}\\
            &=\int{\tan^{n-2}{x}\d\tan{x}}-\int{\tan^{n-2}}{x}\d x\\
            &=\frac{\tan^{n-1}x}{n-1}-\int{\tan^{n-2}{x}\d x}
        \end{aligned}
    }
    有了这个递推公式，我们再结合起始项$\int{\tan{x}\d x}, \int{\tan^2{x}\d x}$的积分便可得到最后结果.
    \formula{\int{\tan{x}\d x}=\int{\frac{\sin{x}}{\cos{x}}\d x}=\int{\frac{\d\cos{x}}{\cos{x}}}=\ln{\abs{\cos{x}}}+C}
    \formula{\int{\tan^2{x}\d x=\int{\qty(\sec^2{x}-1)\d x}=\tan{x}-x+C}}
    后面的递推过程就不展示了，也就是一步步代入找规律，奇偶分开.
    
    结论展示一下，不必记忆，会递推即可.
    \formula{\int{\tan^{2m+1}{x}\d x}=\sum_{r=0}^{m-1}\qty(-1)^{r+2}\frac{\tan^{2m-2r}{x}}{2m-2r}+\qty(-1)^m\ln{\abs{\sec{x}}}+C}
    \formula{\int{\tan^{2m}{x}\d x}=\sum_{r=0}^{m-1}\qty(-1)^r\frac{\tan^{2m-2r-1}{x}}{2m-2r-1}+\qty(-1)^mx+C}
}%在花括号内输入分析